% Part: normal-modal-logic
% Chapter: tableaux
% Section: introduction

\documentclass[../../../include/open-logic-section]{subfiles}

\begin{document}

\olfileid{nml}{tab}{int}

\olsection{प्रस्तावना}

!!^{tableau}s म्हणजे !!{signed
  formula}s चे (खाली शाखा फुटणारे) काही वृक्ष.
चिन्हांकित सूत्र म्हणजे सत्यतामूल्याचे चिन्ह
($\True$ किंवा $\False$) आणि !!a{sentence}
यांची जोडी:
\[
\sFmla{\True}{!A} \text{किंवा} \sFmla{\False}{!A}.
\]
!!^a{tableau} ची सुरुवात काही \emph{गृहीतकां}पासून
होते. पुढील प्रत्येक !!{signed formula} हे अनुमानाचा
एखादा नियम लावून मिळते. काही नियम वृक्षाच्या
टोकावर एक किंवा अधिक !!{signed formula}s जोडतात;
इतर नियम दोन नवी टोके जोडून दोन शाखा तयार करतात.
नियम लावल्यावर मिळणाऱ्या !!{signed formula}s मधील
!!{formula}, नियम ज्या !!{signed formula} ला लावला
त्यातील सूत्रापेक्षा कमी गुंतागुंतीचे असते.
एखाद्या शाखेत $\sFmla{\True}{!A}$ आणि
$\sFmla{\False}{!A}$ ही दोन्ही असतील, तर ती शाखा
\emph{बंद} आहे असे म्हणतो. !!a{tableau} मधील प्रत्येक
शाखा बंद असेल, तर संपूर्ण !!{tableau} बंद असतो.
बंद !!{tableau} ही अशी !!a{derivation} असते की
तिच्यावरून त्या !!{tableau} ची सुरुवात ज्या
!!{signed formula}s ने केली तो संच अपूरणीय
आहे हे दिसते. यावरून $\Proves$ संबंधाची व्याख्या
करता येते: $\Gamma \Proves !A$ तेव्हा आणि तेव्हाच
$\Gamma_0 = \{!B_1,
\dots, !B_n\} \subseteq \Gamma$ असा एखादा सांत
संच असतो की पुढील गृहीतकांसाठी बंद !!{tableau}
मिळतो:
\[
\{\sFmla{\False}{!A}, \sFmla{\True}{!B_1}, \dots, \sFmla{\True}{!B_n}\}.
\]

मोडल तर्कशास्त्रांसाठी !!{signed
formula} ही संकल्पना विस्तारित करावी लागते,
तसेच~\iftag{prvBox}{$\Box$\iftag{prvDiamond}{ आणि
    $\Diamond$}}{$\Diamond$} यांसाठी नियम जोडावे लागतात.
चिन्हाबरोबर ($\True$ किंवा $\False$) मोडल
!!{tableau}s मधील !!{formula}s ना
\emph{पूर्वचिन्हे}~$\sigma$ ही असतात.
ही पूर्वचिन्हे धन पूर्णांकांच्या रिकाम्या नसलेल्या
क्रमिका असतात; म्हणजे $\sigma \in (\PosInt)^* \setminus
\{\emptyseq\}$. अशी पूर्वचिन्हे लिहिताना भोवतालचे
$\tuple{\ }$ वगळतो आणि वेगवेगळ्या !!{element}s
मध्ये $,$ ऐवजी $.$ वापरतो.
$\sigma$ हे पूर्वचिन्ह असेल, तर $\sigma.n$ म्हणजे
$\sigma \concat \tuple{n}$; उदाहरणार्थ
$\sigma = 1.2.1$ असेल, तर $\sigma.3$ म्हणजे
$1.2.1.3$. उदाहरणार्थ,
\[
\sFmla{\True}{\Box !A \lif !A}[1.2]
\]
हे \emph{पूर्वचिन्हयुक्त !!{signed formula}}
(किंवा संक्षेपाने \emph{पूर्वचिन्हयुक्त !!{formula}}) आहे.

सहज अर्थाने, !!a{tableau} च्या एखाद्या शाखेवरील
!!{formula}s पूर्ण करू शकेल अशा प्रतिमानातील
जगाचे नाव पूर्वचिन्ह देते. आणि $\sigma$ हे
एखाद्या जगाचे नाव असेल, तर $\sigma.n$ हे
म्हणजे~$\sigma$ ने दर्शवलेल्या जगापासून प्राप्य
असलेल्या जगाचे नाव देते.

\end{document}
